
\subsection{热力学定律}
\begin{tabularx}{\columnwidth}{XX}
  热力学第一定律    & $\Delta U=Q+W$                           \\
  等容过程       & $W=0$                                    \\
  等压过程气体对外做功 & $W_{\text{气}}=p\Delta V$                 \\
  压强线性变化做功   & $W_{\text{气}}=\frac{p_1+p_2}{2}\Delta V$ \\
  等温过程       & $\Delta U=0$                             \\
  绝热过程       & $Q=0$                                    \\
\end{tabularx}

\textbf{第一定律}

第一定律给出内能变化与做功、热传递之间的关系。等容过程中系统不对外做功；等压过程中可用$W=p\Delta V$计算体积功；等温过程中理想气体内能不变；绝热过程中系统与外界无热交换。

\textbf{自由膨胀}

自由膨胀过程中气体失去外界压强约束，通常视为不做功。

\textbf{第二定律}

热量不能自发地从低温物体传向高温物体；不可能从单一热源吸收热量并全部转化为有用功。

\textbf{第三定律}

绝对零度无法达到。

\textbf{第零定律}

当系统之间停止热交换时，它们达到热平衡，温度相同。

\begin{formulanotebox}
  使用热力学第一定律时需统一符号约定，尤其是外界对系统做功和系统对外做功的正负号，不同教材约定可能不同。本文公式
  \[
    \Delta U=Q+W
  \]
  中，$W$表示外界对气体做功。如果改用气体对外做功$W_{\text{气}}$，则
  \[
    \Delta U=Q-W_{\text{气}}.
  \]
  气缸问题中更常先求气体对外做功，因为体积增大时$W_{\text{气}}>0$，直观上就是气体推动活塞做功。
\end{formulanotebox}

\begin{keypointbox}
  \textbf{热力学三条线}：第一定律管“数量守恒”，第二定律管“方向约束”，第三定律限定“绝对零度不可达”。
\end{keypointbox}

\begin{casetypebox}[title={\strut\textcolor{white}{\bfseries 常见题型：自由膨胀}}]
  自由膨胀过程中气体没有推动外界边界做功，通常取
  \[
    W_{\text{气}}=0.
  \]
  若容器绝热，则
  \[
    Q=0.
  \]
  由热力学第一定律可得
  \[
    \Delta U=0.
  \]
  对理想气体来说，内能只与温度有关，因此自由膨胀且绝热时温度不变。注意自由膨胀不是等压膨胀，不能用$p\Delta V$计算做功。
\end{casetypebox}

\begin{casetypebox}[title={\strut\textcolor{white}{\bfseries 常见题型：气缸等压变化做功}}]
  气缸活塞缓慢移动且外界压强、活塞重力等外部条件不变时，气体常做等压变化。设气体压强为$p$，体积从$V_1$变为$V_2$，则气体对外做功为
  \[
    W_{\text{气}}=p(V_2-V_1).
  \]
  若气体膨胀，$V_2>V_1$，气体对外做正功；若气体被压缩，$V_2<V_1$，气体对外做负功。\\
  例如竖直气缸中活塞质量为$M$，横截面积为$S$，外界大气压为$p_0$，气体在活塞下方。活塞缓慢平衡时
  \[
    p=p_0+\frac{Mg}{S}.
  \]
  所以气体从体积$V_1$膨胀到$V_2$时
  \[
    W_{\text{气}}=\left(p_0+\frac{Mg}{S}\right)(V_2-V_1).
  \]
  再由
  \[
    \Delta U=Q-W_{\text{气}}
  \]
  求热量或内能变化。这个题型对应气体实验定律中的自由活塞等压过程。
\end{casetypebox}

\begin{casetypebox}[title={\strut\textcolor{white}{\bfseries 常见题型：压强均匀变化做功}}]
  如果气体压强随体积线性变化，例如弹簧活塞气缸中
  \[
    p=p_0+\frac{kx}{S},
  \]
  而体积变化为
  \[
    \Delta V=Sx,
  \]
  则$p-V$图像是一条直线。气体对外做功等于$p-V$图像下方面积：
  \[
    W_{\text{气}}=\int_{V_1}^{V_2}p\,dV.
  \]
  当压强从$p_1$均匀变化到$p_2$时，
  \[
    W_{\text{气}}=\frac{p_1+p_2}{2}(V_2-V_1).
  \]
  弹簧活塞中也可以写成外界大气部分和弹簧部分做功之和：
  \[
    W_{\text{气}}=p_0\Delta V+\frac12kx^2.
  \]
  其中$p_0\Delta V$对应推开大气，$\frac12kx^2$对应克服弹簧弹力。\\
  这类题的关键是：压强不是恒定时不能直接用$p\Delta V$，必须用平均压强乘体积变化，或者使用$p-V$图像面积。
\end{casetypebox}

\begin{casetypebox}[title={\strut\textcolor{white}{\bfseries 常见题型：卡口气缸的分段做功}}]
  气缸卡口问题要分阶段计算做功。活塞被卡口挡住时体积不变，所以
  \[
    W_{\text{气}}=0.
  \]
  此阶段即使气体压强随温度升高而增大，也没有体积变化，因此没有体积功。\\
  活塞离开卡口或尚未碰到卡口时，若缓慢平衡且外部条件不变，则为等压过程：
  \[
    W_{\text{气}}=p\Delta V.
  \]
  因此整个过程的做功应按阶段相加：
  \[
    W_{\text{气,总}}=W_{\text{等容段}}+W_{\text{等压段}}.
  \]
  例如先等容升温到卡口临界，再等压膨胀，则
  \[
    W_{\text{气,总}}=0+p(V_2-V_k).
  \]
  最后再用
  \[
    \Delta U=Q-W_{\text{气,总}}
  \]
  处理内能和吸放热。这个题型和气体实验定律中的卡口问题完全对应：先判断体积是否变，再决定有没有做功。
\end{casetypebox}

\begin{casetypebox}[title={\strut\textcolor{white}{\bfseries 常见题型：不同截面积双活塞做功}}]
  不同截面积双活塞由固定杆连接时，气体体积变化为
  \[
    \Delta V=(S_2-S_1)x.
  \]
  若没有额外外力且忽略摩擦，前面气体实验定律中可得气体压强保持$p=p_0$，所以气体对外做功为
  \[
    W_{\text{气}}=p_0(S_2-S_1)x.
  \]
  如果活塞连杆整体还受到恒定外力$F$，平衡时
  \[
    p=p_0+\frac{F}{S_2-S_1}.
  \]
  于是
  \[
    W_{\text{气}}=\left(p_0+\frac{F}{S_2-S_1}\right)(S_2-S_1)x.
  \]
  也就是
  \[
    W_{\text{气}}=p_0(S_2-S_1)x+Fx.
  \]
  这说明气体做功分成两部分：一部分用于排开外界大气，另一部分用于克服外力推动连杆。
\end{casetypebox}
